\chapter{Protocolo de anotação orientada à empresa}
\label{ap:protocolo}

\section{Finalidade e versão do protocolo}

Este apêndice apresenta as regras da tarefa de anotação de notícias em relação à Sabesp. O objetivo é identificar a orientação do acontecimento para a empresa-alvo, conforme o texto disponível, e não avaliar o tom emocional da reportagem ou prever a reação posterior da ação SBSP3. Os resultados apresentados no Capítulo~\ref{cap:resultados} utilizam cem notícias anotadas pelo autor no protocolo orientado à Sabesp, em outubro de 2026. As regras descritas aqui também servem de base para a etapa independente planejada.

\section{Material apresentado ao anotador}

Cada item deverá apresentar identificador estável, título, corpo da notícia, empresa-alvo e campos destinados ao registro do julgamento e de sua justificativa. A versão do texto deve ser preservada, para que eventuais alterações posteriores na página jornalística não modifiquem o material efetivamente avaliado. O anotador deverá utilizar o conteúdo apresentado e não consultar as previsões do modelo, as respostas de outros participantes ou os retornos posteriores da ação durante o julgamento independente.

\section{Pergunta de anotação e categorias}

A pergunta orientadora é: \emph{considerando o título e o corpo da notícia, o acontecimento relatado contém informação favorável, desfavorável ou neutra para a Sabesp?} O julgamento deve estar sustentado pelo conteúdo do texto, sem presumir que todo fato considerado positivo para a companhia provoque valorização de suas ações.

\begin{itemize}
    \item \textbf{Positiva:} o conteúdo apresenta informação favorável para a Sabesp, com fundamento identificável na matéria.
    \item \textbf{Negativa:} o conteúdo apresenta informação desfavorável para a Sabesp, também sustentada no texto.
    \item \textbf{Neutra:} há conteúdo relativo à Sabesp, mas não se identifica orientação favorável ou desfavorável predominante segundo o protocolo.
    \item \textbf{Não avaliável:} o material é insuficiente, tangencial ou inadequado para um julgamento confiável; essa categoria deverá ser registrada separadamente na rodada futura, não convertida automaticamente em neutra.
\end{itemize}

A categoria ``não avaliável'' não foi coletada sistematicamente como classe separada nas cem notícias da rodada atual do autor. Portanto, a introdução dessa categoria na etapa futura exige versionamento do protocolo e cuidado ao comparar métricas produzidas por regras diferentes.

\section{Categorias de foco e casos difíceis}

O tipo de matéria é registrado independentemente do sentimento:
\begin{itemize}
    \item \textbf{Focal:} a Sabesp é o assunto principal ou recebe desenvolvimento substantivo;
    \item \textbf{Agenda ou multiempresa:} a Sabesp aparece junto de outras empresas ou temas gerais de mercado;
    \item \textbf{Tangencial:} a menção não contém informação suficiente para uma avaliação confiável da empresa.
\end{itemize}

Em matérias multiempresa, a polaridade deve se referir apenas à Sabesp. A categoria de foco não determina o sentimento. Quando uma notícia apresenta efeitos mistos ou indiretos, o anotador deverá registrar a evidência e, se necessário, a dificuldade de julgamento, em lugar de atribuir uma classe apenas pela presença de uma palavra positiva ou negativa.

\section{Campos de registro e consolidação da referência}

\begin{table}[htbp]
\centering
\caption{Campos de registro para a anotação orientada à Sabesp.}
\label{tab:campos-anotacao}
\begin{tabularx}{\textwidth}{@{}P{0.32\textwidth}L@{}}
\toprule
Campo & Finalidade \\
\midrule
Identificador da notícia & Relacionar julgamento e versão do corpus \\
Identificação do anotador & Preservar a autoria dos julgamentos independentes \\
Versão do protocolo & Registrar quais regras foram utilizadas \\
Classe de polaridade & Positiva, negativa, neutra ou não avaliável (etapa futura) \\
Tipo de matéria & Focal, agenda/multiempresa ou tangencial \\
Evidência textual & Registrar o trecho que sustenta o julgamento \\
Observação de ambiguidade & Explicitar casos de efeitos mistos ou insuficiência informacional \\
\bottomrule
\end{tabularx}
\vspace{0.4em}
{\footnotesize Fonte: elaboração própria.}
\end{table}

As anotações individuais serão preservadas antes de qualquer referência consolidada. A concordância será calculada somente sobre itens efetivamente avaliados por participantes comparáveis. Havendo divergências, sua resolução será documentada por procedimento de adjudicação, isto é, uma decisão posterior registrada sem apagar os julgamentos independentes. O número de participantes, os itens compartilhados e as regras de consolidação deverão constar da descrição da nova avaliação.

\chapter[Ficha técnica da execução temporal (outubro de 2026)]{Ficha técnica da execução temporal\\(outubro de 2026)}
\label{ap:rastreabilidade}

\section{Síntese do estudo de caso}

Este apêndice registra os principais números do inventário de notícias e as condições da campanha temporal exploratória apresentada no Capítulo~\ref{cap:resultados}. A Tabela~\ref{tab:corpus-p02} descreve a consolidação das notícias do inventário P02, e a Tabela~\ref{tab:rastreabilidade} reúne o universo analisado, os parâmetros e as medidas reportadas. As contagens são provenientes dos registros existentes na campanha, e não de nova execução realizada para esta reescrita.

\begin{table}[htbp]
\centering
\caption{Composição do inventário P02 utilizado na campanha temporal.}
\label{tab:corpus-p02}
\begin{tabularx}{\textwidth}{@{}LrL@{}}
\toprule
Registro & Quantidade & Evidência de execução \\
\midrule
Linhas brutas no inventário & 1.254 & \texttt{selection\_report.json} \\
Notícias canônicas únicas & 1.253 & \texttt{selection\_report.json} \\
Duplicatas mapeadas & 1 & \texttt{selection\_report.json} \\
Exclusões explícitas & 0 & \texttt{selection\_report.json} \\
Notícias com probabilidades válidas, entrada contextual & 1.253 & \texttt{result\_registry.json} \\
Veículos distintos & 5 & \texttt{coverage\_daily.csv} \\
\bottomrule
\end{tabularx}
\vspace{0.4em}
{\footnotesize Fonte: inventário P02 e registros da campanha de outubro de 2026. Distribuição das notícias únicas por veículo: Money Times, 360; InfoMoney, 328; Valor Econômico, 251; G1 Economia, 203; Exame, 111.}
\end{table}

As quantidades da Tabela~\ref{tab:corpus-p02} descrevem o inventário preservado para a execução, não a cobertura completa de todas as notícias publicadas sobre a Sabesp no período. A ausência de um artigo no inventário não prova que nenhum veículo publicou conteúdo relevante naquele dia.

\begin{table}[htbp]
\centering
\caption{Parâmetros e resultados registrados na campanha exploratória ITI (Sabesp/SBSP3).}
\label{tab:rastreabilidade}
\begin{tabularx}{\textwidth}{@{}P{0.34\textwidth}L@{}}
\toprule
Item & Registro \\
\midrule
Empresa e ativo & Sabesp; SBSP3 \\
Período informacional & Maio de 2022 a abril de 2024 (P02) \\
Entrada do classificador para o ITI & Recorte de sentenças sobre a Sabesp, com uma vizinha de cada lado; limite 512 tokens \\
Meia-vida principal & $H=7$ dias civis; sensibilidades $H=2$ e $H=14$ \\
Aquecimento & 70 dias civis, equivalentes a dez semanas iniciais excluídas do contraste \\
Painel semanal alinhado & 104 semanas, chave W-FRI \\
Comparação principal & 94 semanas; $\Delta\rho=+0{,}040$; IC $[-0{,}127;\,0{,}241]$; $p=0{,}512$ \\
Média semanal versus ITI-H7 & 94 semanas; $\Delta\rho=-0{,}062$; IC $[-0{,}165;\,0{,}060]$; $p=0{,}340$; exploratório fora de Holm \\
Correlação marginal ITI-H7 em painel completo & 104 semanas; $\rho\simeq-0{,}145$; inclui aquecimento \\
Bootstrap & 10.000 reamostragens, blocos de quatro semanas; valores $p$ exploratórios \\
Reserva de notícias & Regra \texttt{calendar\_last\_12\_weeks}; 199 matérias; corte em 06/02/2024 \\
Subpainel da reserva & 11 semanas alinhadas, exploração temporal já inspecionada \\
Contraste na reserva & $\Delta\rho=+0{,}186$; IC $[-0{,}057;\,0{,}670]$; $p\simeq0{,}111$; não confirmatório \\
Avaliação humana independente & Pendente; anotação do autor: 100 notícias \\
\bottomrule
\end{tabularx}
\vspace{0.4em}
{\footnotesize Fonte: registros da campanha de outubro de 2026. Intervalos apresentados são bootstrap nominais; não inferir confirmação por seus limites isolados.}
\end{table}

A correlação marginal calculada sobre 104 semanas não deve ser confundida com o contraste principal de 94 semanas, que exclui o aquecimento. Da mesma forma, a reserva de aproximadamente doze semanas de calendário de notícias gerou onze observações semanais alinhadas; o subpainel já foi examinado e não constitui período confirmatório independente.

Para que outra pessoa possa reproduzir exatamente os cálculos, cada execução precisa identificar a versão do corpus, os filtros e a deduplicação, o checkpoint do modelo, o mapeamento de classes, a representação textual, o parâmetro de meia-vida, os preços ajustados efetivamente carregados, as regras temporais e o ambiente de software. A disponibilidade apenas do código e das tabelas finais não substitui os dados e os manifestos associados à execução.

\chapter{Especificação matemática complementar do ITI}
\label{ap:iti-math}

Este apêndice complementa as equações da Seção~\ref{sec:formulacao-iti} com uma derivação curta da meia-vida, a política de ausências e um exemplo numérico. Os valores ilustrativos não são resultados da avaliação temporal da Sabesp.

\section{Derivação da meia-vida}

Na ausência de novas entradas, $X_{c,d}=0$, a recorrência do ITI fica $ITI_{c,d}^{(H)}=\lambda_H ITI_{c,d-1}^{(H)}$. Após $k$ dias consecutivos elegíveis sem entrada, a contribuição remanescente do estado inicial será
\begin{equation}
    ITI_{c,d+k}^{(H)}=\lambda_H^k ITI_{c,d}^{(H)}.
\end{equation}
A definição de meia-vida exige que, para $k=H$, o fator seja $1/2$. Logo, $\lambda_H^H=1/2$, de onde $H\log\lambda_H=-\log 2$ e, finalmente,
\begin{equation}
    \lambda_H=2^{-1/H}.
\end{equation}
A expressão descreve a persistência matemática do indicador, sem afirmar que o efeito econômico de uma notícia se reduza à metade nesse mesmo período.

\section{Política de ausências (resumo)}

\begin{itemize}
    \item \textbf{Dia elegível com notícias:} $X_{c,d}$ é a média simples dos escores das matérias do dia.
    \item \textbf{Dia elegível sem notícia observada:} $X_{c,d}=0$, e o estado do índice decai. Esse zero não corresponde à classificação de uma notícia como neutra.
    \item \textbf{Dia fora da cobertura delimitada:} não se cria sinal para avaliação temporal; ausência não é imputada como silêncio informacional.
\end{itemize}

Dentro da janela P02, a elegibilidade é definida pelo intervalo inventariado, não por comprovação individual de coleta bem-sucedida em todos os dias. A documentação dessa cobertura é uma pendência metodológica, e a interpretação dos zeros precisa reconhecer essa limitação.

\section{Esperança matemática do escore}

Se as classes são codificadas por $Y\in\{-1,0,+1\}$, o escore $s=p^+-p^-$ corresponde à média dos valores de classe ponderada pela distribuição prevista pelo modelo:
\begin{equation}
    \operatorname{E}_{\mathrm{modelo}}[Y] = (+1)p^++(0)p^0+(-1)p^-=p^+-p^-.
\end{equation}
Essa identidade é algébrica e não constitui demonstração de calibração das probabilidades, de acerto na classificação ou de impacto econômico do texto.

\section{Exemplo ilustrativo complementar (não experimental)}

Considere uma notícia para a qual $p^+=0{,}60$, $p^0=0{,}30$ e $p^-=0{,}10$. Seu escore é $s=0{,}50$. Em um dia com três notícias cujos escores sejam $0{,}50$, $-0{,}20$ e $0{,}10$, a entrada diária será
\begin{equation}
    X=\frac{0{,}50-0{,}20+0{,}10}{3}\simeq0{,}1333.
\end{equation}
Partindo de $ITI_0=0$ e adotando $H=7$ (portanto $\lambda_7\simeq0{,}9057$), o primeiro dia produz $ITI_1\simeq(1-0{,}9057)\times0{,}1333\simeq0{,}0126$. Se o dia seguinte for elegível e não contiver notícia observada, $X_2=0$, o Soft será zero e o ITI passará a aproximadamente $0{,}9057\times0{,}0126\simeq0{,}0114$.

\section{Verificações numéricas}

A suíte de testes do projeto contempla casos controlados de meia-vida e decaimento, construção da grade diária, alinhamento de retornos semanais e procedimentos estatísticos. A existência desses testes é parte da documentação da implementação, mas não significa que todos os dados originais ou todos os resultados da campanha tenham sido reproduzidos de modo independente nesta revisão editorial. A reexecução integral depende dos artefatos e configurações identificados no Apêndice~\ref{ap:rastreabilidade}.
